\chapter{Conclusion}
\label{chap:conclusion}
\section{Summary of Research}
\label{sec:summary_of_research}
This thesis focused on the challenge of automating Q\&A functionality of the NamazApp e-learning service by providing timely, trustworthy, and specific answers. The main goal was to design, build, and evaluate an Automatic Learner Support (ALS) system. This system needed to answer simple, "standardized" questions on its own using verified information, while also determining more "complex" questions to assist an expert answering them. The aim was to make the expert's job easier and get users answers faster.
To do this, we used a Retrieval Augmented Generation (RAG) approach. As detailed in Chapter \ref{chap:met}, the ALS system is an LLM-based RAG pipeline that handles active user chats by: (1) using a Large Language Model (LLM) to find unanswered questions and salams; (2) getting relevant information from a Knowledge Base, which includes NamazApp's own Previously Answered Questions (PAQs) and a collection of IslamQA fatwas from trusted sources; (3) using an LLM to create answers based on this retrieved information; and (4) deciding if an answer should be sent automatically or it should be verified by an expert.
\section{Key Findings and Contributions}
\label{sec:key_findings_and_contributions}
We tested the ALS system using 98 real user conversations from NamazApp (see Chapter \ref{chap:eval} for details) and found several important things.
\begin{enumerate}
\item \textbf{Reliable Automatic Answers:} The system achieved 82.54\% precision. This means most of the answers sent directly to users were correct, according to our expert. This supports our main design idea of focusing on reliability by using only PAQs for automatic replies.
\item \textbf{Key Factors for Performance:} The tests showed how important the quality and amount of information in the PAQ database are for good automatic answers. Also, how well the system first extracts the user's questions from the chat really affects whether the final answer is suitable for the situation.
\end{enumerate}
The main contributions of this thesis are:
\begin{itemize}
\item The design and build of a new RAG-based LLM system (ALS) made specifically to give reliable answers to questions about Islam in the NamazApp.
\item A practical way to tell the difference between questions the system can answer on its own ("standardized") and those an expert needs to see ("complex"), focusing on high accuracy by only sending PAQ-based answers automatically.
\item Real-world testing that showed the ALS system can work well, with numbers to back up its reliability and how much it can automate.
\item Learning about the challenges of building RAG systems for specific, sensitive topics, especially how the quality of the information, the instructions given to the LLM, and the rules for sending answers all work together.
\end{itemize}
\section{Limitations of the Study}
\label{sec:conc_limitations}
Even though the results are positive, this project has some limitations, which we covered in Chapter \ref{chap:eval}. The main ones include:
\begin{itemize}
\item The current PAQ database is relatively small, around 300 questions, which limits the number of different questions the system can confidently answer on its own.
\item Although we used real user data, a larger test set could give us even more information.
\end{itemize}
These limitations point to clear ways we can improve things in the future.
\section{Future Research Directions}
\label{sec:conc_future_work}
Based on what we have done, there are several interesting ways this research could go next:
\begin{itemize}
\item Keep expanding and improving the PAQ database.
\item Make the system better at understanding user questions, perhaps by adding an ability to clarify questions.
\item Find more advanced ways to decide when to send an answer, while making sure we keep precision high.
\item Do more tests to compare different LLMs, embedding models, or even other ways of building the system (like fine-tuning versus RAG).
\end{itemize}
\section{Concluding Remarks}
\label{sec:conc_remarks}
This thesis showed that it is possible and valuable to build a reliable LLM-based system to answer questions about Islam in the NamazApp. Our ALS system is efficient and trustworthy, which is really important for a religious application. The system we built already helps by taking some work off the expert and giving users faster answers to common questions.
Our findings highlight how important good quality information sources and clear decision-making rules are when building AI assistants for specific areas. There are still challenges such as needing more data in our database and getting even better at understanding what users mean. Future work on growing our information sources and improving the system will definitely lead to even better and more trusted tools for learning, especially in religious education.